\clearpage
\begin{proofbox}
\textbf{\textcolor{CProof}{方法一：最快证明法（考场两步速写）}}
\begin{center}\begin{tikzpicture}[x=1cm,y=1cm,>=stealth,font=\small]
\path[use as bounding box] (-5.4,-1.3) rectangle (5.4,4.3);
\node[font=\small\bfseries] at (0,4) {速证：两条连线对应两个斜率分母};
\draw[->,black!55] (-4.5,0)--(4.6,0) node[right] {$x$};
\draw[->,black!55] (0,-.6)--(0,3.35) node[above] {$y$};
\draw[blue!65!black,thick,domain=-.5:3.1,samples=70,variable=\t] plot ({sqrt(4+.8*\t*\t)},\t);
\draw[blue!65!black,thick,domain=-.5:3.1,samples=70,variable=\t] plot ({-sqrt(4+.8*\t*\t)},\t);
\coordinate (A1) at (-2,0); \coordinate (A2) at (2,0);
\coordinate (P) at (3,2.5); \coordinate (H) at (3,0);
\node[below left] at (0,0) {$O$};
\draw[orange!85!black,very thick] (A1)--(P);
\draw[teal!80!black,very thick] (A2)--(P);
\fill[black] (P) circle (2pt) node[right=8pt] {$P$};
\fill (A1) circle (2pt) node[below right=4pt] {$A_1$};
\fill (A2) circle (2pt) node[below left=4pt] {$A_2$};

\node[orange!85!black,fill=BGProof,inner sep=3pt] at (-1.35,2.3) {$k_1=\dfrac{y}{x+a}$};
\node[teal!80!black,fill=BGProof,inner sep=3pt] at (4,1.5) {$k_2=\dfrac{y}{x-a}$};
\node at (0,-1) {$P\in C:\quad y^2=\dfrac{b^2}{a^2}(x^2-a^2)$};
\end{tikzpicture}
\end{center}


\textbf{识别信号：}题目出现\(P\)到两顶点的斜率积；只求定值，不必求\(P\)坐标。
设\(P(x,y)\in C\setminus\{A_1,A_2\}\)，则
\[
\boxed{k_{PA_1}k_{PA_2}=\frac{y^2}{(x+a)(x-a)}
 =\frac{y^2}{x^2-a^2}}.
\]
而\(\displaystyle\frac{x^2}{a^2}-\frac{y^2}{b^2}=1\)立刻给出
\[
\boxed{\frac{y^2}{x^2-a^2}=\frac{b^2}{a^2}}.
\]
故\(k_{PA_1}k_{PA_2}=\dfrac{b^2}{a^2}=e^2-1\)。\quad\(\square\)

\textbf{为什么最快：}仅写“\textbf{斜率相乘出现平方差}\(\to\)\textbf{曲线方程整体替换}”，
无需解出\(x,y\)，也无需联立直线。全程不展开多余式子。

\medskip
\end{proofbox}
\clearpage
\begin{proofbox}
\textbf{\textcolor{CProof}{方法二：最低门槛证明法（两个直角三角形）}}
\begin{center}\begin{tikzpicture}[x=1cm,y=1cm,>=stealth,font=\small]
\path[use as bounding box] (-5.4,-1.3) rectangle (5.4,4.3);
\node[font=\small\bfseries] at (0,4) {作垂足：共同的高，两个不同的底};
\draw[->,black!55] (-4.5,0)--(4.6,0) node[right] {$x$};
\draw[->,black!55] (0,-.6)--(0,3.35) node[above] {$y$};
\draw[blue!65!black,thick,domain=-.5:3.1,samples=70,variable=\t] plot ({sqrt(4+.8*\t*\t)},\t);
\draw[blue!65!black,thick,domain=-.5:3.1,samples=70,variable=\t] plot ({-sqrt(4+.8*\t*\t)},\t);
\coordinate (A1) at (-2,0); \coordinate (A2) at (2,0);
\coordinate (P) at (3,2.5); \coordinate (H) at (3,0);
\node[below left] at (0,0) {$O$};
\draw[orange!85!black,very thick] (A1)--(P);
\draw[teal!80!black,very thick] (A2)--(P);
\fill[black] (P) circle (2pt) node[right=8pt] {$P$};
\fill (A1) circle (2pt) node[below right=4pt] {$A_1$};
\fill (A2) circle (2pt) node[below left=4pt] {$A_2$};

\fill[orange!50,opacity=.15] (A1)--(H)--(P)--cycle;
\fill[teal!60,opacity=.2] (A2)--(H)--(P)--cycle;
\draw[dashed,black!70] (P)--(H) node[midway,right=5pt] {$|y|$};
\draw (2.8,0)--(2.8,.2)--(3,.2);
\node[below right] at (H) {$H$};
\draw[<->,orange!85!black] (-2,-.95)--(3,-.95) node[midway,fill=BGProof] {$|x+a|$};
\draw[<->,teal!80!black] (2,3.15)--(3,3.15) node[midway,above] {$|x-a|$};
\draw[dotted,black!40] (2,.2)--(2,3.15) (3,2.7)--(3,3.15);

\end{tikzpicture}
\end{center}


取\(H=(x,0)\)为\(P(x,y)\)向实轴作垂线的垂足。
因\(P\)在双曲线上且不在顶点，\(|x|>a\)，\(|y|>0\)。
观察直角三角形\(\triangle PA_1H\)、\(\triangle PA_2H\)：
\[
 |k_{PA_1}|=\frac{|PH|}{|A_1H|}=\frac{|y|}{|x+a|},\qquad
 |k_{PA_2}|=\frac{|PH|}{|A_2H|}=\frac{|y|}{|x-a|}.
\]
\textbf{第一眼能看出：}两条斜率同号，所以斜率积等于绝对值之积。
因此
\[
 k_{PA_1}k_{PA_2}
 =\frac{|y|^2}{|x+a|\,|x-a|}
 =\frac{y^2}{x^2-a^2}.
\]
\textbf{第二眼使用定义方程：}
\(\frac{y^2}{b^2}=\frac{x^2-a^2}{a^2}\)，
即\(\frac{y^2}{x^2-a^2}=\frac{b^2}{a^2}\)。证毕。

\textbf{为什么更容易理解：}从“高÷底”的初中直角三角形斜率出发，
靠共同的高\(|PH|\)和两段水平距离发现平方差；
\textbf{先理解符号，再计算比值}，左右两支均适用。

\medskip
\end{proofbox}

\clearpage
\begin{proofbox}
\textbf{\textcolor{CProof}{方法三：圆幂定理证明（初中知识的几何升级）}}
\begin{center}% Exact model: a=2, b=sqrt(5), P=(3,5/2), H=(3,0).
% T=(4/3,-2sqrt(5)/3), so OT^2=4, OT dot HT=0, HT^2=5.
\begin{tikzpicture}[x=1.1cm,y=1.1cm,>=stealth,font=\small]
\draw[rounded corners=5pt,black!25] (-3.7,-3.65) rectangle (5.1,4.5);
\node[font=\bfseries] at (.7,3.9) {圆幂桥梁：共同高 $PH$ 与切线长 $HT$};
\coordinate (O) at (0,0);
\coordinate (A1) at (-2,0); \coordinate (A2) at (2,0);
\coordinate (H) at (3,0); \coordinate (P) at (3,2.5);
\coordinate (T) at ({4/3},{-2*sqrt(5)/3});
\draw[->,black!50] (-3.15,0)--(4.55,0) node[right] {$x$};
\draw[->,black!50] (0,-2.25)--(0,3.25) node[above] {$y$};
\draw[blue!65!black,thick,domain=0:3.15,samples=80,variable=\t]
 plot ({sqrt(4+.8*\t*\t)},\t);
\draw[blue!65!black,thick,domain=0:3.15,samples=80,variable=\t]
 plot ({-sqrt(4+.8*\t*\t)},\t);
\draw[violet!75!black,dashed,thick,domain=0:360,samples=120,variable=\t]
 plot ({2*cos(\t)},{2*sin(\t)});
\draw[orange!85!black,thick] (A1)--(P);
\draw[teal!80!black,thick] (A2)--(P);
\draw[black!75,very thick] (P)--(H);
\node[right=7pt] at (3,1.25) {$PH=|y|$};
\draw[red!65!black,very thick] (H)--(T);
\node[red!65!black] at (2.45,-1.25) {$HT$};
\draw[violet!75!black,thick] (O)--(T);
\node[below left=2pt] at (.65,-.72) {$a$};
\draw[black!60] ($(T)!.18cm!(O)$)--($(T)!.18cm!(O)+(T)!.18cm!(H)-(T)$)--($(T)!.18cm!(H)$);
\draw (2.82,0)--(2.82,.18)--(3,.18);
\fill (O) circle (1.6pt) node[above left=3pt] {$O$};
\fill (A1) circle (1.8pt) node[below left=5pt] {$A_1$};
\fill (A2) circle (1.8pt) node[below left=5pt] {$A_2$};
\fill (H) circle (1.8pt) node[below right=5pt] {$H(x,0)$};
\fill (P) circle (1.8pt) node[right=8pt] {$P(x,y)$};
\fill (T) circle (1.8pt) node[below right=7pt] {$T$};
\node[violet!75!black] at (-2.75,-1.7) {辅助圆};
\draw[dotted,black!35] (-2,-.3)--(-2,-2.6) (3,-.4)--(3,-3.1) (2,-.4)--(2,-3.1);
\draw[<->,black!70] (-2,-2.6)--(3,-2.6) node[midway,fill=white,inner sep=3pt] {$HA_1=x+a$};
\draw[<->,black!70] (2,-3.1)--(3,-3.1);
\node[left=5pt] at (2,-3.1) {$HA_2=x-a$};
\end{tikzpicture}
\end{center}
\textbf{为什么想到作圆：}两斜率相乘，分母会出现两个水平距离
$HA_1$ 和 $HA_2$ 的乘积；初中\textbf{切割线定理}专门把它变成一个平方。

\textbf{第一步：圆幂。}
以 $A_1A_2$ 为直径作圆，设 $H(x,0)$ 是 $P$ 在 $x$ 轴上的垂足。
$|x|>a$ 保证 $H$ 在圆外。从 $H$ 作切线 $HT$，得到
\[
 HA_1\cdot HA_2=HT^2=OH^2-a^2=x^2-a^2.
\]
这里 $HA_1,HA_2$ 均是\textbf{非负长度}，不是带符号坐标。

\textbf{第二步：斜率积变为长度比平方。}
两斜率同号，故
\[
 k_1k_2=\frac{PH}{HA_1}\cdot\frac{PH}{HA_2}
 =\boxed{\left(\frac{PH}{HT}\right)^2}.
\]

\textbf{第三步：用曲线方程固定长度比。}
$PH^2=y^2$，$HT^2=x^2-a^2$，由 $P\in C$：
\[
 \frac{PH^2}{b^2}=\frac{HT^2}{a^2}
 \quad\Longrightarrow\quad \frac{PH}{HT}=\frac ba,
 \quad\boxed{k_1k_2=\frac{b^2}{a^2}}.
\]
\textbf{价值：}把双曲线第三定义接到初中圆的\textbf{圆幂定理}；
但最后一步仍使用双曲线方程，不能误称全程纯几何。\quad\(\square\)
\end{proofbox}

\clearpage
\begin{proofbox}
\textbf{\textcolor{CProof}{方法四：常规解析几何证明（规范三步）}}
\begin{center}\begin{tikzpicture}[x=1cm,y=1cm,>=stealth,font=\small]
\path[use as bounding box] (-5.4,-1.3) rectangle (5.4,4.3);
\node[font=\small\bfseries] at (0,4) {规范设点：坐标决定斜率的增量};
\draw[->,black!55] (-4.5,0)--(4.6,0) node[right] {$x$};
\draw[->,black!55] (0,-.6)--(0,3.35) node[above] {$y$};
\draw[blue!65!black,thick,domain=-.5:3.1,samples=70,variable=\t] plot ({sqrt(4+.8*\t*\t)},\t);
\draw[blue!65!black,thick,domain=-.5:3.1,samples=70,variable=\t] plot ({-sqrt(4+.8*\t*\t)},\t);
\coordinate (A1) at (-2,0); \coordinate (A2) at (2,0);
\coordinate (P) at (3,2.5); \coordinate (H) at (3,0);
\node[below left] at (0,0) {$O$};
\draw[orange!85!black,very thick] (A1)--(P);
\draw[teal!80!black,very thick] (A2)--(P);
\fill[black] (P) circle (2pt) node[right=8pt] {$P$};
\fill (A1) circle (2pt) node[below right=4pt] {$A_1$};
\fill (A2) circle (2pt) node[below left=4pt] {$A_2$};

\draw[dashed,black!50] (0,2.5) node[left] {$y_0$}--(P)--(H);
\node[below right] at (H) {$x_0$};
\node[above left] at (A1) {$-a$}; \node[above left] at (A2) {$a$};
\node at (0,-1) {$P=(x_0,y_0),\quad |x_0|>a,\quad y_0\ne0$};
\end{tikzpicture}
\end{center}


\begin{description}[style=nextline,leftmargin=2.8em,itemsep=0.4em]
 \item[\textbf{步骤1（明确定义域并设点）：}]
 设\(P(x_0,y_0)\in C\)，且\(P\ne A_1,A_2\)。
 因\(x_0^2/a^2=1+y_0^2/b^2>1\)，有\(x_0\ne\pm a\)。
 \item[\textbf{步骤2（写斜率）：}]
 \[
 k_{PA_1}=\frac{y_0}{x_0+a},\quad
 k_{PA_2}=\frac{y_0}{x_0-a},\quad
 k_{PA_1}k_{PA_2}=\frac{y_0^2}{x_0^2-a^2}.
 \]
 \item[\textbf{步骤3（代入标准方程）：}]
 \[
 \frac{x_0^2}{a^2}-\frac{y_0^2}{b^2}=1
 \ \Longrightarrow\ 
 y_0^2=\frac{b^2}{a^2}(x_0^2-a^2).
 \]
 故\(k_{PA_1}k_{PA_2}=\frac{b^2}{a^2}\)。
 又由\(c^2=a^2+b^2\)，有
 \(\frac{b^2}{a^2}=\frac{c^2}{a^2}-1=e^2-1\)。
\end{description}

\textbf{为什么需要：}有清楚的条件检查、斜率定义与方程代换，适合书面完整证明。

\medskip
\end{proofbox}
\clearpage
\begin{proofbox}
\textbf{\textcolor{CProof}{方法五：逆向轨迹证明（真正的“第三定义”）}}
\begin{center}\begin{tikzpicture}[x=1cm,y=1cm,>=stealth,font=\small]
\path[use as bounding box] (-5.4,-1.3) rectangle (5.4,4.3);
\node[font=\small\bfseries] at (0,4) {逆向轨迹：曲线保留，两个顶点删除};
\draw[->,black!55] (-4.5,0)--(4.6,0) node[right] {$x$};
\draw[->,black!55] (0,-.6)--(0,3.35) node[above] {$y$};
\draw[blue!65!black,thick,domain=-.5:3.1,samples=70,variable=\t] plot ({sqrt(4+.8*\t*\t)},\t);
\draw[blue!65!black,thick,domain=-.5:3.1,samples=70,variable=\t] plot ({-sqrt(4+.8*\t*\t)},\t);
\coordinate (A1) at (-2,0); \coordinate (A2) at (2,0);
\coordinate (P) at (3,2.5); \coordinate (H) at (3,0);
\node[below left] at (0,0) {$O$};
\draw[orange!85!black,very thick] (A1)--(P);
\draw[teal!80!black,very thick] (A2)--(P);
\fill[black] (P) circle (2pt) node[right=8pt] {$P$};
\draw[blue!65!black,fill=white,line width=1pt] (A1) circle (2pt) node[below right=4pt] {$A_1$};
\draw[blue!65!black,fill=white,line width=1pt] (A2) circle (2pt) node[below left=4pt] {$A_2$};

\node[orange!85!black,fill=BGProof] at (-1.2,1.45) {$k_1$};
\node[teal!80!black,fill=BGProof] at (3.8,1.4) {$k_2$};
\node at (0,-1) {空心点不属于轨迹；图取 $a=2,\ m=5/4$};
\end{tikzpicture}
\end{center}


\textbf{已知：}两固定点\(A_1(-a,0),A_2(a,0)\)（\(a>0\)），
任意点\(P(x,y)\)与它们的连线斜率均存在，且
\(k_{PA_1}k_{PA_2}=m>0\)。
则\(x\ne\pm a\)，于是
\[
 \frac{y^2}{(x+a)(x-a)}=m
 \quad\Longleftrightarrow\quad
 y^2=m(x^2-a^2)
 \quad\Longleftrightarrow\quad
 \boxed{\frac{x^2}{a^2}-\frac{y^2}{a^2m}=1}.
\]
反过来，取此双曲线上\textbf{非顶点}\(P\)，按斜率公式可逐步逆推得斜率积\(m\)。
故轨迹为该双曲线\textbf{去掉两顶点}，不是整条含顶点的双曲线。

\textbf{为什么值得补充：}使“\textbf{曲线上点\(\Rightarrow\)斜率积定值}”升级为
“\textbf{斜率积定值\(\Longleftrightarrow\)双曲线轨迹}”，
遇到动点轨迹题可以倒着用。

\medskip
\end{proofbox}

\clearpage
\begin{proofbox}
\textbf{\textcolor{CProof}{方法六：三角半角参数法（高阶巧证）}}
% [TikZ待绘] C065_half_angle.tex：x=a sec(theta), y=b tan(theta)，显示两斜率及渐近线。
\IfFileExists{assets/tikz/C065_half_angle.tex}{%
 \begin{center}\input{assets/tikz/C065_half_angle.tex}\end{center}
}{}
\textbf{为什么值得用参数：}如果 $P$ 在右支，作
\[
 x=a\sec\theta,\quad y=b\tan\theta,
 \qquad -\frac\pi2<\theta<\frac\pi2,\ \theta\ne0.
\]
因为 $\sec^2\theta-\tan^2\theta=1$，这个参数化\textbf{自动保证 $P$ 在双曲线上}。

\textbf{第一条斜率：}
\[
 k_1=\frac{b\tan\theta}{a(\sec\theta+1)}
 =\frac ba\frac{\sin\theta}{1+\cos\theta}
 =\boxed{\frac ba\tan\frac\theta2}.
\]
\textbf{第二条斜率：}
\[
 k_2=\frac{b\tan\theta}{a(\sec\theta-1)}
 =\frac ba\frac{\sin\theta}{1-\cos\theta}
 =\boxed{\frac ba\cot\frac\theta2}.
\]
于是
\[
 \boxed{k_1k_2=\left(\frac ba\right)^2
 \tan\frac\theta2\cot\frac\theta2=\frac{b^2}{a^2}}.
\]
左支由中心对称 $(x,y)\mapsto(-x,-y)$ 得到；两条\textbf{有序}斜率互换，乘积不变。
\textbf{理解价值：}原本看似无关的两斜率，分别包含\textbf{互为倒数}的半角因子。
\textbf{适用边界：}此法需要三角参数化，不如方法一适合限时考试。\quad\(\square\)
\end{proofbox}

\clearpage
\begin{proofbox}
\textbf{\textcolor{CProof}{纵轴型独立核对（防机械套用）}}
\begin{center}\begin{tikzpicture}[x=1cm,y=1cm,>=stealth,font=\small]
\path[use as bounding box] (-5.4,-3.1) rectangle (5.4,4.1);
\node[font=\bfseries] at (0,3.8) {纵轴型：共同的水平增量，竖直增量不同};
\draw[->,black!55] (-3,0)--(4.7,0) node[right] {$x$};
\draw[->,black!55] (0,-2.7)--(0,3.1) node[left] {$y$};
\draw[blue!65!black,thick,domain=-2.5:3.5,samples=80,variable=\t] plot (\t,{sqrt(1+\t*\t/4)});
\draw[blue!65!black,thick,domain=-2.5:3.5,samples=80,variable=\t] plot (\t,{-sqrt(1+\t*\t/4)});
\coordinate (P) at (3,1.802776);
\draw[orange!85!black,very thick] (0,-1)--(P);
\draw[teal!80!black,very thick] (0,1)--(P);
\fill (0,-1) circle (2pt) node[above left=5pt] {$B_1(0,-a)$};
\fill (0,1) circle (2pt) node[below left=5pt] {$B_2(0,a)$};
\fill (P) circle (2pt) node[above left=5pt] {$P(x,y)$};
\draw[dashed] (P)--(3,-1)--(0,-1);
\draw[dotted,black!45] (0,1)--(4.3,1) (P)--(4.3,1.802776);
\draw[<->,teal!80!black] (4.3,1)--(4.3,1.802776) node[midway,right] {$y-a$};
\node[right=6pt] at (3,.35) {$y+a$};
\draw[<->] (0,-2.4)--(3,-2.4) node[midway,fill=BGProof] {$x$};
\node at (0,-3) {图取 $a=1,\ b=2$；斜率始终为 $\Delta y/\Delta x$};
\end{tikzpicture}
\end{center}

对\(\frac{y^2}{a^2}-\frac{x^2}{b^2}=1\)及
\(B_1(0,-a), B_2(0,a)\)，非顶点处\(x\ne0\)，
\[
 k_{PB_1}k_{PB_2}
 =\frac{(y+a)(y-a)}{x^2}
 =\frac{y^2-a^2}{x^2}
 =\frac{a^2}{b^2}=\frac{1}{e^2-1}.
\]
\end{proofbox}
